\documentclass[main.tex]{subfiles}
\begin{document}
\ifSubfilesClassLoaded{\begin{center}\small\itshape Part of the GnssProcessor Mathematical Reference \textbullet\ \today\end{center}}{}
\section{Integer ambiguity resolution}\label{sec:ar}
% ============================================================
%  Integer ambiguity resolution
%  MLAMBDA over a between-satellite-differenced ambiguity sub-problem
%  (BSSD in single-receiver modes, true DD in baseline) extracted from
%  the float Kalman state, with WL/NL
%  decomposition, partial/subset acceptance, ratio + FFR
%  validation, and constrained pseudo-observation feedback.
%  Impl: driver in src/core/filter/AmbiguityResolver.{hpp,cpp};
%  MLAMBDA numerical kernels + acceptance battery in
%  src/core/filter/IlsCore.{hpp,cpp} (ils:: free functions);
%  network-producer AR seams (both default-OFF) in
%  src/core/filter/NetworkProducerAr.{hpp,cpp} (log-only fixability
%  probe) and src/core/filter/NetworkSdEwlAr.{hpp,cpp} (SD
%  extra-wide-lane integer decision, Route B + A2 datum)
% ============================================================

Integer ambiguity resolution (AR) maps the real-valued, bias-contaminated
phase ambiguities of the float solution to the integer carrier-cycle
counts they represent.  The engine runs
\textsc{MLAMBDA}~\cite{teunissen1995lambda} on a
between-satellite-differenced ambiguity sub-problem extracted from the
filter state and covariance.  The operative quantity is the ambiguity
differenced against a per-constellation pivot $\piv{g}$
(\dcref{sec:diff}{the differencing section of the Foundations
document}).  The single-receiver and baseline modes share the
integer-estimation machinery, but the condition that makes the
differenced ambiguity an integer differs by mode.

The undifferenced ambiguity $N^\sat_{\rx,\sigl}$ of the phase
observation equation (\dcref{eq:obs-phase}{Foundations, observation
equations}) carries a receiver phase bias $b_{\rx,\sigl}$ and a
satellite phase bias $b^\sat_\sigl$.  In single-receiver PPP and
PPP-RTK the pivot difference is a \emph{between-satellite single
difference} $\bsd{\,\cdot\,}$ (\dcref{sec:bssd}{the between-satellite
single-difference section of the Foundations document}).  It cancels
the receiver bias $b_{\rx,\sigl}$, which is common to every tracked
satellite, and leaves the satellite bias.  The residual
$\bsd N^{\sat,\piv{g}}_{\rx,\sigl}$ is therefore integer only after
$b^\sat_\sigl$ is removed by a consistent external phase-bias product
(BIA or SSR, \dcref{sec:ssr}{the SSR-corrections document}).  In
baseline mode the state holds true double differences $\dd{\,\cdot\,}$
(\dcref{sec:dd}{the double-difference section of the Foundations
document}); the additional between-receiver difference also cancels the
satellite bias, so $\dd N$ is integer without an external product.  The
equations below keep the resolver's internal label $\dd{}$ for the
differenced ambiguity, since the algebra is the same in both cases: it
denotes the between-satellite single difference $\bsd{}$ in the
single-receiver case and the full double difference only in baseline
mode.

The implementation separates a driver from a numerical kernel.  The
driver, \texttt{src/core/filter/AmbiguityResolver.cpp}, assembles the
state sub-vector, rescales between metres and cycles
(metre$\leftrightarrow$cycle), looks up the per-pair $\lambda$, forms the
differences, performs the WL/NL decomposition and applies the returned
integer fix.  The numerical core lives in the
\texttt{GnssProcessor::ils} namespace of
\texttt{src/core/filter/IlsCore.\{hpp,cpp\}}: the $LDL^{\Tr}$
factorisation (in the factor order of \eqref{eq:ldlt}), the
$\Zmat$-reduction, the Schnorr--Euchner search, the bootstrap success
rate and the full acceptance battery.  These are pure free functions over
$(\ambfloat\ \text{cycles},\,\ambcov\ \text{cycles}^2)$, with no
dependence on \texttt{StateManager}, \texttt{KalmanFilter} or the
configuration.  The PPP-RTK network producer (\dcref{sec:net}{the
Network mode document}) calls the same \texttt{ils} core under its own
gate set, as a read-only fixability probe only, because the producer is
float-only (\Cref{sec:ar-network}).  The single-receiver and baseline
paths below are configured under the JSON path
\texttt{estimation.ambiguity\_resolution.*}.

\subsection{The integer least-squares problem}\label{sec:ar-ils}

The float solution gives a real-valued ambiguity estimate $\ambfloat$
(in cycles) with covariance $\ambcov$, both extracted from the filter
state $\stateV$ and covariance $\covM$.  Integer ambiguity resolution is
the weighted closest-lattice-point problem

\begin{empheq}[box=\fbox]{equation}
\ambfix
\;=\;
\arg\min_{\bm{a}\,\in\,\mathbb{Z}^{n}}
\;\bigl(\ambfloat-\bm{a}\bigr)^{\Tr}\,
\ambcov^{-1}\,
\bigl(\ambfloat-\bm{a}\bigr),
\label{eq:ils}
\end{empheq}

that is, the minimisation of the squared Mahalanobis distance from
$\ambfloat$ to a lattice point $\bm{a}$ in the metric of $\ambcov^{-1}$.
This integer least-squares (ILS) problem is NP-hard.  \textsc{LAMBDA} and
its \textsc{MLAMBDA} variant make it tractable in three stages, none of
which has a tuning parameter: (i) an $LDL^{\Tr}$ factorisation of
$\ambcov$, (ii) a decorrelating integer $\Zmat$-transform (lattice
reduction), and (iii) a depth-first shrinking-ellipsoid search over the
reduced lattice.  The hat/check convention of the notation block
applies: $\ambfloat$ is the float estimate and $\ambfix$ the
integer-fixed solution.

\paragraph{Float extraction and rescaling to cycles.}
The filter stores per-band ambiguity states in metres
($\lambda_\sigl N$, one per tracked satellite and signal).  The resolver
therefore forms the between-satellite-differenced float and covariance
by explicit differencing against the pivot $\piv{g}$, then rescales by
wavelength.  In baseline mode the state already holds the double
difference, and the differencing step reduces to a direct selection
(\Cref{sec:ar-feedback}).  For a differenced pair $k$ with satellite
index $i(k)$ and pivot index $p(k)$
(\texttt{src/core/filter/AmbiguityResolver.cpp}):
\begin{align}
y^{\dd{}}_k &= \stateV_{\,i(k)} - \stateV_{\,p(k)},
  &&\text{(DD float, m)} \label{eq:dd-float}\\[2pt]
\bm{Q}^{\dd{}}_{kl} &=
  \covM_{i_k i_l} - \covM_{i_k p_l}
  - \covM_{p_k i_l} + \covM_{p_k p_l},
  &&\text{(DD covariance, m}^2\text{)} \label{eq:dd-cov}\\[2pt]
\ambfloat_k = y^{\dd{}}_k / \lambda_k,
\qquad
(\ambcov)_{kl} &= \bm{Q}^{\dd{}}_{kl} \big/ (\lambda_k \lambda_l).
  &&\text{(cycle domain)} \label{eq:m2c}
\end{align}
The rescaling \eqref{eq:m2c} is required.  On metre-domain floats
\textsc{MLAMBDA} would round toward integer-metre values, which have no
carrier-cycle meaning, and the ratio test (\Cref{sec:ar-accept}) would
return $\approx 1$ irrespective of float quality.  The per-pair
wavelength $\lambda_k$ uses the GLONASS FDMA per-satellite frequency
where applicable.  $\bm{Q}^{\dd{}}$ is symmetrised,
$\bm{Q}^{\dd{}} \leftarrow \tfrac12(\bm{Q}^{\dd{}} + \bm{Q}^{\dd{}\Tr})$,
before factorisation.

The resolver elects its own DD reference: in each
$(\text{constellation},\,\text{signal})$ group of AR-eligible
ambiguities, the highest-elevation member is the reference and every
other member is differenced against it.  This reference differs from
the filter's BSSD pivot.  Under BSSD the pivot's ambiguity is frozen at
zero variance, which excludes it from the candidate set
(\Cref{sec:ar-partial}); the resolver therefore differences against a
second reference, and the resulting quantity is a difference of single
differences.  The two elections also use different rules.  The BSSD
pivot is chosen on a hard phase-coverage tier, a convergence/elevation
score and an ephemeris-freshness bar (\dcref{sec:diff}{the differencing
section of the Foundations document}); the AR reference is chosen on
elevation alone (\texttt{AmbiguityResolver::buildDdPairs}).

\subsection{Decorrelation: $LDL^{\Tr}$ factorisation and the integer
$\Zmat$-transform}\label{sec:ar-lambda}

\paragraph{Factorisation.}
\textsc{MLAMBDA} factorises the symmetric positive-definite ambiguity
covariance as
\begin{equation}
\ambcov \;=\; \Lmat^{\Tr}\,\Dmat\,\Lmat,
\qquad
\Lmat \text{ unit lower-triangular},\;\;
\Dmat = \operatorname{diag}(d_1,\dots,d_n),\;\; d_i > 0,
\label{eq:ldlt}
\end{equation}
eliminated in reverse index order, from the last row up (the Chang--Yang
\textsc{mlambda} convention).  The transposed unit
lower-triangular factor stands on the left, as written in
\eqref{eq:ldlt}; the shorthand ``LDL-transpose'' used in this section
names that factorisation, not the reversed factor order.  A
non-positive pivot $d_i \le 0$ signals a covariance that is not positive
definite and aborts the fix.  After reduction, $d_i$ is the
\emph{conditional variance} of the $i$-th decorrelated ambiguity given
the later ones; the success-rate gate of \Cref{sec:ar-accept} uses it
directly.

\paragraph{The $\Zmat$-transform.}
The central device of \textsc{LAMBDA}~\cite{teunissen1995lambda} is a
unimodular integer transform $\Zmat \in \mathbb{Z}^{n\times n}$,
$|\det\Zmat| = 1$, that re-parameterises
\begin{equation}
\bm{z} = \Zmat^{\Tr}\ambfloat,
\qquad
\bm{Q}_{\hat z} = \Zmat^{\Tr}\,\ambcov\,\Zmat,
\label{eq:ztransform}
\end{equation}
so that the transformed ambiguities are nearly decorrelated (small
off-diagonal $\Lmat$ entries) and their conditional variances are
ordered.  The decorrelated float vector is
$\bar{\bm{z}} := \Zmat^{\Tr}\ambfloat$, with components $\bar{z}_k$; it
is the quantity the search of \Cref{sec:ar-lambda} rounds, computed as
\texttt{z\_float = Z.transpose()*a} in \texttt{ils::mlambda}
(\texttt{src/core/filter/IlsCore.cpp}).  Because $\Zmat$ is unimodular,
$(\Zmat^{\Tr})^{-1}$ is also integer and the integer minimiser maps back
exactly, $\ambfix = (\Zmat^{\Tr})^{-1}\check{\bm{z}}$.  Decorrelation
therefore leaves the solution of \eqref{eq:ils} unchanged.  It makes the
search ellipsoid closer to spherical, which reduces the number of
enumerated points by orders of magnitude.  The reduction interleaves two
operations in a sweep from $j = n-2$ downward.

\emph{Integer Gauss transformations} reduce the sub-diagonal $\Lmat$
entries by integer column operations,
\begin{equation}
\mu = \lceil L_{ij} \rfloor,
\qquad
\Lmat_{(:,j)} \mathrel{-}= \mu\,\Lmat_{(:,i)},
\qquad
\Zmat_{(:,j)} \mathrel{-}= \mu\,\Zmat_{(:,i)},
\label{eq:gauss}
\end{equation}
where $\lceil\cdot\rfloor$ is round-to-nearest, keeping every
$|L_{ij}| \le \tfrac12$.

\emph{Permutations} (conditional adjacent swaps) reorder two ambiguities
when the swap reduces the later conditional variance.  With
$\bar{d}_j = d_j + L_{j+1,j}^2\,d_{j+1}$, the swap fires when
\begin{equation}
\bar{d}_j + \varepsilon \;<\; d_{j+1},
\qquad \varepsilon = 10^{-6},
\label{eq:perm-test}
\end{equation}
after which $d_j,\,d_{j+1}$, the affected rows/columns of $\Lmat$ and
the columns of $\Zmat$ are updated by the standard rank-one formulas
($\eta = d_j/\bar{d}_j$, $\lambda = d_{j+1}L_{j+1,j}/\bar{d}_j$,
$d_j \leftarrow \eta\,d_{j+1}$, $d_{j+1} \leftarrow \bar{d}_j$), and the
sweep restarts from $j = n-2$.  The guard $\varepsilon = 10^{-6}$
prevents repeated swaps of near-equal variances.  It is a numerical
safeguard rather than an estimator setting and has no configuration key.

\paragraph{The shrinking-ellipsoid search.}
The reduced problem is solved by depth-first enumeration of the lattice
points inside the ellipsoid $\sum_i (z_i - \hat{z}_i)^2 / d_i \le
\chi^2$, in \mbox{Schnorr--Euchner} order (nearest integer first, then
zig-zagging outward).  Recursing from the last level $k = n-1$ down to
$0$,
\begin{align}
\hat{z}_k &= \bar{z}_k + \sum_{i>k}(z_i - \hat{z}_i)\,L_{ik},
  &&\text{(conditional centre)} \label{eq:cond}\\
z_k = \lceil \hat{z}_k \rfloor,
\qquad
\mathrm{dist}_k &= \mathrm{dist}_{k+1}
  + \frac{(\hat{z}_k - z_k)^2}{d_k}.
  &&\text{(partial distance)} \label{eq:dist}
\end{align}
A complete integer vector that reaches level $0$ with cumulative
distance below the current bound $\chi^2$ is recorded.  Once the
candidate list is full, the bound shrinks to the worst retained
candidate, so the ellipsoid contracts as better candidates are found.
The search (\texttt{ils::search} in
\texttt{src/core/filter/IlsCore.cpp}) retains the $m = 2$ best
candidates for the ratio test, sorted by ascending squared distance,
under a hard iteration cap (\texttt{kMaxIter} $= 10\,000$).  Fewer than
two candidates is a failure.  The outputs are the two smallest objective
values $s_1 \le s_2$ of \eqref{eq:ils} and the best integer vector
$\check{\bm{z}}$, back-transformed via
$\ambfix = (\Zmat^{\Tr})^{-1}\check{\bm{z}}$.

\subsection{Wide-lane / narrow-lane decomposition (PPP
path)}\label{sec:ar-wlnl}

PPP and PPP-RTK AR use a \emph{wide-lane-first} decomposition when
Melbourne--W\"ubbena (MW) resolved wide-lane (WL) integers are
available.  It halves the search dimension and exploits the long WL
wavelength~\cite{laurichesse2009integer}.  The driver attempts the
WL-constrained narrow-lane (NL) fix before any unconstrained
\textsc{LAMBDA}.  When WL integers exist and the constrained fix fails,
the epoch stays float: the driver does not fall through to an
unconstrained search, because that search would include DDs whose WL is
not locked and could return integers that contradict the MW wide lane.
A wrong-WL fix is worse than no fix, so the enforced rule is ``no NL AR
for satellites whose WL is float''
(\texttt{src/core/filter/AmbiguityResolver.cpp}).

\paragraph{WL from Melbourne--W\"ubbena (integer by contract).}
The DD wide-lane integer is not searched.  It is the difference of the
MW detector's single-difference WL integers
(\texttt{src/core/slip/MwSlipDetector}):
\begin{equation}
\dd N_{\mathrm{WL}} = N^{\,\sat}_{\mathrm{WL}} - N^{\,\piv{g}}_{\mathrm{WL}}
  \in \mathbb{Z},
\qquad
\lambda_{\mathrm{WL}} = \frac{\sLight}{f_1 - f_2}.
\label{eq:dd-wl}
\end{equation}

\paragraph{NL search variable.}
For each dual-frequency DD pair the transform $\bm{T}$ forms the float
narrow-lane ambiguity in NL cycles.  The default row is the narrow-lane
combination below.  It is not ionosphere-free: its slant-ionosphere
sensitivity is 0.959 NL cycles per TECU, so the default per-pair cap
\eqref{eq:perpair} of 0.15 cycles implicitly requires the slant
ionosphere to about 0.156 TECU.
\begin{equation}
\hat{n}_{\mathrm{NL}}
= \frac{f_1\,\dd N_1 + f_2\,\dd N_2}{f_1 + f_2}
= \frac{1}{f_1 + f_2}\!\left(
    \frac{f_1}{\lambda_1}\,\dd{}_1 + \frac{f_2}{\lambda_2}\,\dd{}_2
  \right),
\qquad
\lambda_{\mathrm{NL}} = \frac{\sLight}{f_1 + f_2},
\label{eq:nl-float}
\end{equation}
with reduced covariance $\bm{Q}_{\mathrm{NL}} = \bm{T}\,\bm{Q}^{\dd{}}\,
\bm{T}^{\Tr}$.  Only NL-paired DDs enter \textsc{LAMBDA}; unpaired
single-signal DDs are excluded so that they cannot distort the ratio
test.  An optional ionosphere-free search row (configuration key
nl\_iono\_free\_row, default off) targets the same L1 integer with zero
slant-ionosphere sensitivity.  For GPS L1/L2 it multiplies the WL
sensitivity by 8.06 (a 1-cycle WL error moves the L1 integer target by
3.53 cycles instead of 0.438) and the searched variance by about 65.
The equations of this subsection describe the default narrow-lane row.

\paragraph{Integer-recentring shift.}
The float $\hat{n}_{\mathrm{NL}}$ is not integer-valued in general.
Substituting $N_2 = N_1 - \mathrm{WL}$ gives
$\hat{n}_{\mathrm{NL}} = N_1 - \tfrac{f_2}{f_1+f_2}\,\mathrm{WL}$, which
for GPS L1/L2 with $\mathrm{WL} = 1$ lies at $N_1 - 0.4379$.  Neither
neighbouring integer is then materially closer, and the ratio test settles
near $1.6$ regardless of float quality.  The resolver therefore adds the
known-WL offset, which makes the search variable integer in
expectation,
\begin{equation}
\tilde{n}_k
= \hat{n}_{\mathrm{NL},k} + \frac{f_2}{f_1 + f_2}\,\dd N_{\mathrm{WL},k}
\;\approx\; N_{1,k} \;\in\; \mathbb{Z} + \text{noise},
\label{eq:nl-shift}
\end{equation}
runs \textsc{MLAMBDA} on $\tilde{\bm{n}}$ (a constant per-row shift
leaves $\bm{Q}_{\mathrm{NL}}$ unchanged), then back-shifts
$\check{n}_{\mathrm{NL}} = \check{N}_1 -
\tfrac{f_2}{f_1+f_2}\,\mathrm{WL}$.

\paragraph{Back-transform to per-signal DDs.}
The fixed integers map back to the two per-signal DDs in metres,
\begin{equation}
\dd{}_1^{\mathrm{fix}}
  = \Bigl(\check{n}_{\mathrm{NL}}
    + \tfrac{f_2}{f_1+f_2}\,\dd N_{\mathrm{WL}}\Bigr)\lambda_1,
\qquad
\dd{}_2^{\mathrm{fix}}
  = \Bigl(\check{n}_{\mathrm{NL}}
    - \tfrac{f_1}{f_1+f_2}\,\dd N_{\mathrm{WL}}\Bigr)\lambda_2,
\label{eq:nl-back}
\end{equation}
and re-enter the filter through the constrained update of
\Cref{sec:ar-feedback}.

\paragraph{State-basis variants.}
When \texttt{estimation.state\_ambiguity\_basis} $=$ \texttt{nl\_wl} the
state holds NL/WL ambiguities directly, and there are no per-band DD
pairs to enumerate.  In this basis only the MW-resolved WL integers are
applied, as constraint rows; a full NL \textsc{LAMBDA} search with the
WL pre-constrained is not implemented.  The complementary WL-only mode
(\texttt{estimation.ambiguity\_resolution.wl\_only}, in the
\texttt{per\_band} basis) applies the cycle-domain constraint
$\dd{}_{1}/\lambda_1 - \dd{}_{2}/\lambda_2 = \dd N_{\mathrm{WL}}$ as one
pseudo-observation per pair, without an NL search.  When MW re-locks at
a new integer, the covariance of the affected single-difference states
is inflated so that the state moves to the new integer within one epoch.

Both WL-constraint paths (this one and the \texttt{nl\_wl}-basis
constraint) take their pseudo-observation noise from a separate
setting, $\bm{R} = \sigma^2_{\mathrm{wl}}\,\bm{I}$ with
$\sigma^2_{\mathrm{wl}} = \texttt{wl\_only\_pseudo\_var}$ ($10^{-4}$
cycles$^2$, in the cycle domain), rather than the metre-domain
$\sigma^2_{\mathrm{hold}}$ of the NL/DD feedback in \eqref{eq:pseudo}.
Both paths apply three acceptance gates to the MW integers.  The first
is the per-arc fractional-residual point test (\texttt{wl\_max\_frac},
default $0.40$; on both arcs in the \texttt{per\_band} \texttt{wl\_only}
path, on the single arc in the NL/WL basis).  The second is a bootstrap
success-rate gate (key wl\_only\_ffr, default on): the product of
per-arc rounding success rates, taken over distinct arcs, must reach the
threshold min\_success\_rate of \eqref{eq:ffr}, and the weakest
arc is dropped until it does.  The third is an innovation screen (key
wl\_only\_max\_innov\_cyc, default 0.5 cycle) that rejects a constraint
whose asserted WL integer differs from the predicted DD wide lane by
more than that amount.  The cross-frequency parity check runs only in
the \texttt{nl\_wl} basis (\texttt{applyMwWlConstraintsNlWl\_} in
\texttt{src/core/filter/AmbiguityResolver.cpp}).  For a
triple-frequency satellite it rejects all three WL DDs when the
telescoping canonical cycle counts are inconsistent,
$\dd{}_{ab} - \dd{}_{ac} + \dd{}_{bc} \ne 0$.  The \texttt{per\_band}
\texttt{wl\_only} path has no cross-frequency consistency check.

\subsection{Partial / subset ambiguity resolution}\label{sec:ar-partial}

Partial AR operates at two levels: a static candidate gate that decides
which ambiguities enter the integer problem, and an optional
backward-elimination search that looks for a fixable subset when the
full set fails validation.

\paragraph{Upstream candidate gating (always on).}
Subset selection begins when the DD pairs are built
(\texttt{src/core/filter/AmbiguityResolver.cpp}).  A satellite/signal
ambiguity enters the integer set only if it passes every filter:
\begin{itemize}[leftmargin=*, itemsep=0pt]
\item \textbf{Initialised}: the state entry is non-zero (seeded by
      data), and \textbf{non-frozen}: its variance is non-zero, which
      excludes a BSSD-frozen pivot.
\item \textbf{Elevation}: $\theta^\sat_\rx \ge \theta_{\mathrm{AR}}$
      (\texttt{elevation\_mask}).  In the single-receiver path the
      elevation comes from the epoch's observations; an ambiguity whose
      satellite has no observation is given elevation 0 and is removed
      by any positive mask (the consequence for fix-and-hold is stated
      in \Cref{sec:ar-feedback}).
\item \textbf{Resolvable}: the ambiguity is not flagged unresolvable
      for lack of bias or antenna (BIA, ATX) coverage.  This applies to
      the single-receiver path; in baseline DD the analogous filter is
      the \texttt{exclude\_keys} set of non-integer-capable DDs from
      \texttt{MeasurementModel}.
\item \textbf{Constellation}: $g(\sat) \in$
      \texttt{active\_systems}, when that list is non-empty.
\item \textbf{Pairable}: the $(\text{constellation},\,\text{signal})$
      group holds $\ge 2$ satellites, so a pivot difference exists.
\item \textbf{Mature} (\texttt{min\_arc\_info}, default $0$ = off):
      both constituent arcs of a candidate (the higher-frequency arcs of
      the satellite and of the reference) have accumulated at least
      $\sum_k \sin^2\theta^\sat_{\rx,k}$ of elevation-weighted code
      information since arc birth, in epoch-equivalents at zenith.
      Maturity is measured in information rather than epochs: under the
      elevation stochastic model an arc rising off the mask and one at
      zenith accumulate information at very different rates, so two
      arcs of equal age are not equally determined.  The counter
      restarts on both ways an arc can end: the key leaving the state
      vector, and a refloat in place.  A refloat keeps the key and would
      otherwise let the new arc inherit the pre-slip age
      (\texttt{AmbiguityResolver::refreshArcAges\_}, keyed on
      \texttt{StateManager::isAmbiguityPendingReseed}).  The gate applies
      only to the WL-constrained NL candidate set of \Cref{sec:ar-wlnl},
      not to the baseline DD, WL-only or unconstrained full-set paths.
      It guards against wrong fixes, since one immature candidate can
      fail the joint ratio test of \eqref{eq:ratio} for the whole set.
      It does not speed convergence: it admits a candidate only once the
      candidate is well determined, which is when the float solution
      gains least from a fix.
\end{itemize}
The constellation filter keeps a system in the float solution while
excluding it from the candidate set.  It serves a constellation whose
residual model bias would lower the ratio test.

\paragraph{Dynamic backward-elimination subset
(\texttt{enable\_partial\_ar}, default off).}
With the candidate set fixed, the default is all-or-nothing: the full
set is searched and either passes all gates of \Cref{sec:ar-accept} or
the epoch stays float.  When \texttt{enable\_partial\_ar} is set, a
failure of the full-set battery starts a greedy backward elimination
(\texttt{ils::selectFixableSubset} in
\texttt{src/core/filter/IlsCore.cpp}, called through the driver-side
\texttt{selectFixableSubset\_} configuration adapter) before the epoch
falls back to float.  A failure of \textsc{MLAMBDA} itself on the full
set (a covariance that is not positive definite, or fewer than two
candidates) does not start the elimination.  Starting from the full
candidate set, each iteration re-runs the complete \textsc{MLAMBDA}
pipeline of \Crefrange{sec:ar-ils}{sec:ar-lambda} on the current subset
and re-tests the same acceptance battery: the ratio lower and upper
bounds \eqref{eq:ratio}, the bootstrap success rate \eqref{eq:ffr}, and
the per-entry fractional cap \eqref{eq:perpair}.  On any failure it
drops one entry and retries, down to a single ambiguity:
\begin{itemize}[leftmargin=*, itemsep=0pt]
\item if the \emph{fractional cap} is the failing gate, the entry with
      the largest $|\ambfloat_k - \ambfix_k|$ (nearest a half-cycle
      boundary) is removed;
\item otherwise (including a failure of \textsc{MLAMBDA} on the subset)
      the entry with the largest formal float variance $(\ambcov)_{kk}$
      is removed: the least precise ambiguity, whose removal is expected
      to improve the bootstrap rate and the search conditioning most.
\end{itemize}
The accepted subset is always re-validated by the full battery on its
own covariance sub-block.  The procedure can therefore only fix fewer
ambiguities than the full set, never accept a set that fails the
battery: a poor elimination order costs fixes, not integrity.  With the
option off, the path is bit-identical to the all-or-nothing behaviour.
The same routine serves the per-band SD, baseline DD and WL-constrained
NL paths and operates only in the cycle domain, independent of the
decorrelation.  Each caller maps the surviving indices back to its own
state rows and applies them through the constrained update of
\Cref{sec:ar-feedback}.

\subsection{Validation: ratio test and fixed-failure-rate
gate}\label{sec:ar-accept}

A fix is committed only if it passes all three of the following gates.
The code evaluates them in order (ratio lower bound, ratio upper bound,
bootstrap success rate, per-pair cap) and stops at the first failure.
The ratio and fixed-failure-rate (FFR) gates bound different things:
the ratio bounds the discrimination between the two best candidates, the
FFR gate the absolute success probability.

\paragraph{Ratio test (discrimination).}
\begin{empheq}[box=\fbox]{equation}
\mathrm{ratio} = \frac{s_2}{s_1} \;\ge\; \tau_{\mathrm{ratio}},
\qquad s_1 \le s_2 \text{ the two smallest ILS objectives},
\label{eq:ratio}
\end{empheq}
with $s_1, s_2$ from the search of \Cref{sec:ar-lambda} and threshold
$\tau_{\mathrm{ratio}}$ (\texttt{ratio\_threshold}).  An upper-bound
rejection, $\mathrm{ratio} > \tau_{\mathrm{ratio}}^{\max}$
(\texttt{ratio\_upper\_bound}), catches a degenerate covariance
direction: a near-zero $s_1$, typically left by an earlier wrong-integer
hold, produces a very large but meaningless ratio.  A non-positive squared
distance for either candidate rejects the fix before the ratio is formed.

\paragraph{Fixed-failure-rate gate (bootstrap success rate).}
The gate uses a lower bound on the ILS success probability that follows
from the problem geometry,
\begin{empheq}[box=\fbox]{equation}
P_{s,\mathrm{boot}}
= \prod_{i=1}^{n}
  \operatorname{erf}\!\left(\frac{1}{2\sqrt{2\,d_i}}\right)
\;\ge\; P_s^{\min},
\label{eq:ffr}
\end{empheq}
where the $d_i$ are the post-reduction $LDL^{\Tr}$ conditional variances
of \eqref{eq:ldlt} and $P_s^{\min}$ is \texttt{min\_success\_rate}.  Each
factor is $\Pr\bigl(|\mathcal{N}(0,d_i)| < \tfrac12\bigr)$, the
probability that sequential bootstrapped rounding lands on the correct
integer in that dimension.  The product is Teunissen's lower bound on the
ILS success rate.  The acceptance threshold is therefore tied to the
problem geometry rather than to a fixed ratio alone.

\paragraph{Per-pair residual backstop.}
The joint ratio test can pass on a borderline set in which one pair lies
near a half-cycle boundary.  The fix is therefore rejected if the
fractional residual of any single DD exceeds a cap,
\begin{equation}
\max_k \bigl|\ambfloat_k - \ambfix_k\bigr|
  \;\le\; \kappa_{\mathrm{frac}}
\quad\text{(cycles)},
\label{eq:perpair}
\end{equation}
with $\kappa_{\mathrm{frac}}$ = \texttt{nl\_max\_per\_pair\_frac}.

\paragraph{Convergence guards.}
Before any search the resolver skips AR while the float solution is
immature: during its first \texttt{min\_epochs} invocations (a global
counter of resolver calls, not a per-candidate age), and while the
position standard deviation $\sqrt{\covM_{00} + \covM_{11} +
\covM_{22}}$ exceeds \texttt{max\_position\_sigma\_m}.  A marginal fix
accepted before the covariance settles can lock onto wrong integers and
corrupt every later epoch.

\subsection{Feedback: constrained pseudo-observation
update}\label{sec:ar-feedback}

Accepted integers are fed back through a finite-variance pseudo-observation
rather than the classical perfect back-substitution (the $R = 0$
conditioning form)
\begin{equation}
\stateV_{\mathrm{fix}}
  = \stateV - \bm{Q}_{\bm{a}\bm{a}}\,\bm{Q}^{\dd{}-1}
    \bigl(y^{\dd{}} - \ambfix\bigr),
\qquad
\covM_{\mathrm{fix}}
  = \covM - \bm{Q}_{\bm{a}\bm{a}}\,\bm{Q}^{\dd{}-1}
    \bm{Q}_{\bm{a}\bm{a}}^{\Tr},
\label{eq:backsub}
\end{equation}
The update is a \emph{Joseph-form pseudo-observation}
(\texttt{src/core/filter/AmbiguityResolver.cpp}; kernel
\texttt{KalmanFilter::josephPseudoUpdate} in
\texttt{include/filter/KalmanFilter.hpp}).  The reason is structural.
\eqref{eq:backsub} drives the fixed ambiguity rows of $\covM$ to
$\approx 0$ variance, so the next epoch's \textsc{LAMBDA} sees a
near-singular $\bm{Q}^{\dd{}}$ and returns large, meaningless ratios
that keep tripping the upper-bound guard of \eqref{eq:ratio}.  A
finite-variance pseudo-observation is intended to keep $\covM$
well-conditioned, so that later epochs can re-confirm the integer
(converging toward it) or detect a slip through its residual.

This intent is not yet validated.  Reducing $\sigma^2_{\mathrm{hold}}$
(over the range 1e-3 to 1e-7, in square metres) reduced the fix count and
made the kinematic solution diverge, which is the signature of a
constraint that drives the covariance singular rather than tightening
it.  The resolver therefore carries a read-only conditioning probe
(\texttt{qbCondition} / \texttt{rebuildQbCycles} in
\texttt{src/core/filter/AmbiguityResolver.cpp}) that eigen-decomposes
the cycle-domain DD covariance before every search and after every fix,
and logs the conditioning on every fix and every rejection.  The
mechanism described below is the design intent, not a validated result.

For each accepted DD pair $k$ the constraint row, innovation and noise
are
\begin{empheq}[box=\fbox]{equation}
\jacH_{(k,\,i)} = +1,\quad
\jacH_{(k,\,p)} = -1,
\qquad
v_k = \lambda_k\,\ambfix_k - \bigl(\stateV_i - \stateV_p\bigr),
\qquad
\bm{R} = \sigma^2_{\mathrm{hold}}\,\bm{I},
\label{eq:pseudo}
\end{empheq}
followed by
$(\stateV_{\mathrm{fix}}, \covM_{\mathrm{fix}}) =
\textsc{JosephUpdate}(\stateV, \covM, \jacH, \bm{v}, \bm{R})$, with
$\sigma^2_{\mathrm{hold}}$ = \texttt{hold\_variance}.  The DD/baseline
path uses a single-column selector $\jacH_{(k,\,\mathrm{idx})} = +1$,
because the state value already is the double difference.  The same
update applies a partial-AR subset, with fewer rows.  The updated state
and covariance replace the filter's on every epoch with an accepted fix,
in every mode other than 0.  The resolver does not distinguish modes 1
(continuous) and 2 (instantaneous); only mode 0 (float) and mode 3
(fix-and-hold) change its behaviour.  The $R = 0$ form
\eqref{eq:backsub} is not implemented; it appears in the code only in
comments, as the rejected alternative.

Fix-and-hold mode (\texttt{estimation.ambiguity\_resolution.mode}
$= 3$) applies a second pseudo-observation update on each epoch with an
accepted fix, holding the differenced ambiguities at their fixed values
with the same hold variance.  In the baseline path
(\texttt{AmbiguityResolver::holdAmbiguitiesDd}) the hold covers every DD-keyed
ambiguity in the state, including keys excluded from the fix.  In the
single-receiver path (\texttt{AmbiguityResolver::holdAmbiguities}, called by both
the post-processing and the real-time engines) the hold
is inert under the default configuration.  It rebuilds the DD pairs
with an empty observation list, so every ambiguity is assigned elevation
0 and the AR elevation mask (default \SI{20}{\degree}) removes every
pair.  PPP fix-and-hold therefore holds nothing unless the mask is 0 or
negative; mode 3 then behaves like modes 1 and 2.  Under the NL/WL
state basis the hold is a no-op by design.

\subsection{Network-producer AR (both mechanisms default
off)}\label{sec:ar-network}

The PPP-RTK network producer of \dcref{sec:net}{the Network mode
document} is float-only.  Its per-epoch seam sequence contains two
integer mechanisms, both default-off; with the defaults, neither changes
the emitted satellite-bias product.  They are documented here because
their algebra is settled, not because they are in use.

\paragraph{Producer AR: a log-only fixability probe.}
Under \texttt{network.ambiguity\_resolution.enable} (default
\texttt{false}) the seam \texttt{filter/NetworkProducerAr} reconstructs
each station's float ambiguities from the reduced solution
(\texttt{NetworkEstimator::reconstructStation}) and forms an
independent between-station double-difference basis: a star topology
per $(\text{constellation},\,\text{signal})$, restricted to GPS and
Galileo because GLONASS FDMA leaves a per-slot receiver term that
survives the difference.  It then builds the exact DD covariance from
\texttt{NetworkEstimator::jointConditionalCovariance}, rescales it to
cycles by \eqref{eq:m2c}, and calls \texttt{ils::ilsFix}.  The result is
logged and discarded: nothing is re-pinned or re-solved, and the run is
byte-identical whether the block executes or not.  The probe is
meaningful because the between-station DD ambiguity
\begin{equation}
\dd N^{\sat,\piv{g}}_{\rx_1\rx_2,\sigl}
\;=\;
\bsd N^{\sat,\piv{g}}_{\rx_1,\sigl} - \bsd N^{\sat,\piv{g}}_{\rx_2,\sigl}
\label{eq:net-dd-invariant}
\end{equation}
is invariant under the producer's choice of phase-bias S-basis: the
satellite bias $b^\sat_\sigl$ cancels in the between-receiver
difference exactly as it does in the baseline DD of
\dcref{sec:dd}{the double-difference section of the Foundations
document}.  The probe therefore reports the network's intrinsic integer
recoverability rather than a property of the datum the producer holds.
Its convergence gauge is the median DD formal standard deviation in
cycles.  That value must fall well below the half-cycle scale before
\eqref{eq:ils} is reachable, and it is a better progress indicator than
the fix count.  Two configuration notes apply.
\texttt{ils\_max\_dim} (default $0$ = uncapped) exists because the
search does not return on the full-network basis at operational station
counts; a positive value restricts the search to the
lowest-formal-$\sigma$ subset, in the manner of \Cref{sec:ar-partial}.
\texttt{wl\_cascade} is parsed and echoed in the load-time log but has
no consumer in the code, so it is inert.

\paragraph{SD extra-wide-lane integers: Route B and the A2 datum.}
The second mechanism (\texttt{filter/NetworkSdEwlAr}, gated on
\texttt{network.sd\_ewl.fix\_enable}, default \texttt{false}) decides an
integer for the between-satellite single-differenced extra-wide-lane
ambiguity.  The motivation is a rank argument.  The float
Melbourne--W\"ubbena row that the enclosing SD-EWL layer folds is a
linear combination of observables already in the epoch's normal
equations, so it adds no rank; the only new information the row can
supply is integrality.  Rounding the SD statistic directly is circular,
because it requires the satellite bias difference
$b^\sat_{\mathrm{EWL}} - b^{\piv{g}}_{\mathrm{EWL}}$ that the network
estimates.  Route B avoids the circularity by differencing the SD MW
statistic between stations, which gives the same double difference as
\eqref{eq:net-dd-invariant}:
\begin{empheq}[box=\fbox]{equation}
\bigl(\mwc^\sat_{\rx_1} - \mwc^{\piv{g}}_{\rx_1}\bigr)
- \bigl(\mwc^\sat_{\rx_2} - \mwc^{\piv{g}}_{\rx_2}\bigr)
\;=\;
\lamewl\,\dd N^{\sat,\piv{g}}_{\rx_1\rx_2,\mathrm{EWL}}
\;+\; \dd{\varepsilon^{\mwc}},
\label{eq:sdewl-routeb}
\end{empheq}
in which the satellite bias pair cancels exactly: the satellite pair is
the same on both sides, and each station's receiver term has already
cancelled inside its own single difference.  What remains is an integer
plus averaged MW noise at the extra-wide-lane wavelength.  An integer
decided this way uses no a-priori satellite bias and cannot be corrupted
by a wrong one, which is why Route B is preferred over rounding against
a bias product.  The cancellation in \eqref{eq:sdewl-routeb} has a
precondition that belongs to the algebra, not to the implementation:
$b_{\mathrm{EWL}}$ is a per-signal combination, so both stations' arcs
must use the same two tracking channels.  Otherwise a per-signal bias
difference is folded in as though it were an integer.  The seam
therefore elects a network-common channel pair per satellite by station
consensus and admits only stations that hold that pair.

Equation \eqref{eq:sdewl-routeb} yields a double difference, while the
constrained row needs a single-differenced integer per (station,
satellite).  The A2 datum supplies one.  Within each connected component
a reference station $\rx_0$ is elected, $\bsd n^{\sat,\piv{g}}_{\rx_0} :=
0$ is defined, and the rounded DD fixes are propagated outward along a
spanning tree; the link satellite absorbs the per-station free integer.
The asserted quantity is then
\begin{equation}
\bsd n^{\sat,\piv{g}}_{\rx}
\;=\;
\bsd N^{\sat,\piv{g}}_{\rx,\mathrm{EWL}}
- \bsd N^{\sat,\piv{g}}_{\rx_0,\mathrm{EWL}},
\label{eq:sdewl-a2}
\end{equation}
the true SD integer minus a constant per (satellite pair, component),
which is the classical wide-lane fractional-cycle-bias convention modulo
$\lamewl$~\cite{laurichesse2009integer,geng2013triple}.  The constant is
absorbed into the estimated bias the network reports.  Two consequences
of \eqref{eq:sdewl-a2} matter for the design.  First, the per-station
free integer, and with it the spanning tree, cancels at the point of
use, so a change of spanning tree between epochs is not a defect.
Second, $\rx_0$'s own integers are absorbed into the constant, so the
choice of $\rx_0$ is a datum decision rather than a measurement.

Acceptance is the scalar case of the battery in \Cref{sec:ar-accept}.  A
rounded DD is admitted on its fractional residual, as in
\eqref{eq:perpair}, and on the $n = 1$ instance of the bootstrap bound
\eqref{eq:ffr}, each with its own threshold (network.sd\_ewl.fix\_max\_frac\_cyc,
default 0.15 cycle, and network.sd\_ewl.fix\_min\_ffr, default 0.999):
\begin{equation}
P_{s,\mathrm{round}}
= \operatorname{erf}\!\left(\frac{1}{2\sqrt{2}\,\sigma}\right),
\qquad \sigma \text{ in cycles},
\label{eq:sdewl-ffr}
\end{equation}
so a single rounding and the $n$-dimensional search share the same
failure-rate semantics.

The mechanism is subject to the standing rule that a stiff, time-linked
quantity needs a wired event response.  The fixed row writes a
$\lamewl$-quantised value into satellite bias states with very small
process noise, so an integer that changes between epochs would inject a
step that those states take hours to relax out of.  The wired response
is withhold-on-disagreement.  When this epoch's integer differs from the
previous one on a continuous arc, the fix for that arc is withheld
rather than repaired, because on a disagreement epoch neither value is
known to be right.  An arc that restarted (slip, gap, re-acquisition)
replaces the stored integer without counting as a disagreement, because
its integer has legitimately changed.  The stored integer is used only
to veto, never to supply a value; this keeps the mechanism short of a
carried zero-process-noise hold, which is deliberately not built.

\paragraph{Status.}  Both mechanisms are default-off, and neither has a
graded, datum-invariant field result.  The loader also disarms
\texttt{network.sd\_ewl.fix\_enable}, with a warning, unless the float
row layer (\texttt{network.sd\_ewl.enable}) is armed.  That layer is
itself inert without \texttt{network.phase\_leveling} and an MW-class
slip detector, because it reads the arc accumulator they build.  The
emitted network product is float.

\subsection{Configuration summary}\label{sec:ar-config}

All tunable thresholds live under the JSON path
\texttt{estimation.ambiguity\_resolution.*}, parsed in
\texttt{src/core/config/GnssConfigLoad.cpp} with defaults in
\texttt{include/filter/AmbiguityResolver.hpp}.  The estimator theory of
\Crefrange{sec:ar-ils}{sec:ar-feedback} has no settings, since changing
it would change the mathematics rather than a policy.  In the
mode row of the table, modes 1 and 2 behave identically in the
code (\Cref{sec:ar-feedback}).  The table omits three WL acceptance keys
described in \Cref{sec:ar-wlnl} (wl\_only\_ffr, default on;
wl\_only\_min\_arc\_count, default 10; wl\_only\_max\_innov\_cyc, default
0.5) and the search-row option nl\_iono\_free\_row (default off).  The
two network-producer mechanisms of \Cref{sec:ar-network} are configured
separately, under \texttt{network.ambiguity\_resolution.*} and
\texttt{network.sd\_ewl.*} (parsed in
\texttt{src/core/config/GnssConfigLoadNetwork.cpp}); both master switches
default to \texttt{false} and are omitted from the table.

\begin{center}
\footnotesize
\begin{tabular}{@{}lllp{0.32\linewidth}@{}}
\toprule
Symbol & Config key (\texttt{...ambiguity\_resolution.}) & Default & Role \\
\midrule
                          & \texttt{mode}                 & \texttt{0}      & 0 float / 1 continuous / 2 instantaneous / 3 fix-and-hold \\
$\tau_{\mathrm{ratio}}$   & \texttt{ratio\_threshold}     & \texttt{3.0}    & ratio-test acceptance \eqref{eq:ratio} \\
$\tau_{\mathrm{ratio}}^{\max}$ & \texttt{ratio\_upper\_bound} & \texttt{1e4} & degenerate-covariance reject \eqref{eq:ratio} \\
$P_s^{\min}$              & \texttt{min\_success\_rate}   & \texttt{0.999}  & FFR gate \eqref{eq:ffr} \\
$\kappa_{\mathrm{frac}}$  & \texttt{nl\_max\_per\_pair\_frac} & \texttt{0.15} & per-pair residual cap \eqref{eq:perpair} \\
$\theta_{\mathrm{AR}}$    & \texttt{elevation\_mask}  & \texttt{20\textdegree} & candidate elevation mask \\
$\sigma^2_{\mathrm{hold}}$ & \texttt{hold\_variance}      & \texttt{0.001} m$^2$ & pseudo-obs $\bm{R}$ \eqref{eq:pseudo} \\
$n_{\min}^{\text{ep}}$    & \texttt{min\_epochs}          & \texttt{10}     & convergence guard \\
$\sigma_{\mathrm{pos}}^{\max}$ & \texttt{max\_position\_sigma\_m} & \texttt{0.05} m & convergence guard \\
$\textstyle\sum\sin^2\!\theta_{\min}$ & \texttt{min\_arc\_info} & \texttt{0.0} (off) & per-candidate arc-maturity gate, WL/NL path only (\Cref{sec:ar-partial}) \\
                          & \texttt{wl\_only}             & \texttt{false}  & MW-WL-only, no NL search \\
$\sigma^2_{\mathrm{wl}}$  & \texttt{wl\_only\_pseudo\_var} & \texttt{1e-4} cyc$^2$ & WL-constraint pseudo-obs $\bm{R}$ (\Cref{sec:ar-wlnl}) \\
                          & \texttt{wl\_max\_frac}        & \texttt{0.40}   & per-arc WL fractional-residual gate (\Cref{sec:ar-wlnl}) \\
                          & \texttt{enable\_partial\_ar}  & \texttt{false}  & backward-elimination subset fix (\Cref{sec:ar-partial}) \\
                          & \texttt{active\_systems}  & \texttt{[]} (all) & constellation subset (\Cref{sec:ar-partial}) \\
                          & \texttt{state\_ambiguity\_basis}\,$^{\dagger}$ & \texttt{per\_band} & \texttt{per\_band} vs \texttt{nl\_wl} dispatch \\
\bottomrule
\end{tabular}

{\footnotesize $^{\dagger}$\,Path is \texttt{estimation.state\_ambiguity\_basis}, one level above the \texttt{ambiguity\_resolution} block.}
\end{center}

\ifSubfilesClassLoaded{\printbibliography}{}
\end{document}
